\subsection{Fibonacci Search}

% ---------------------------------------------------------------
% 1. Explicación gráfica
% ---------------------------------------------------------------
\begin{frame}{Fibonacci Search: Explicación Gráfica (1/3)}
    \begin{block}{Concepto principal}
        \footnotesize
        Dada $f$ \textbf{unimodal} en $[a,b]$ y un presupuesto fijo de $n$ evaluaciones,
        se consultan dos puntos interiores $c<d$: por unimodalidad el mínimo queda del
        lado del menor valor y el resto se descarta. Las posiciones usan razones de
        Fibonacci $\rho=F_n/F_{n+1}$, que maximizan la reducción garantizada.
    \end{block}

    \begin{center}
    \begin{tikzpicture}[x=0.82cm, y=1cm, font=\scriptsize]
        % Intervalo original con las dos consultas
        \draw[CelesteBrillante, thick] (0,0) -- (8,0);
        \foreach \x/\lb in {0/a, 3/c, 5/d, 8/b}{
            \draw[CelesteBrillante, thick] (\x,-0.13) -- (\x,0.13);
            \node[above, text=GrisTexto] at (\x,0.13) {$\lb$};
        }
        % Caso 1
        \draw[NaranjaBase, very thick] (0,-0.75) -- (5,-0.75);
        \node[right, text=GrisTexto] at (8.25,-0.75) {$[a,d]$ \ si $f(c)<f(d)$};
        % Caso 2
        \draw[NaranjaBase, very thick] (3,-1.4) -- (8,-1.4);
        \node[right, text=GrisTexto] at (8.25,-1.4) {$[c,b]$ \ si $f(c)\ge f(d)$};
    \end{tikzpicture}
    \end{center}

    \scriptsize
    \begin{itemize}
        \item Con $n$ consultas el intervalo final mide $(b-a)/F_{n+1}$: la
              \textbf{máxima reducción garantizada} sin usar derivadas.
        \item El punto que sobrevive se reutiliza: \textbf{1 evaluación nueva} por iteración.
    \end{itemize}
\end{frame}

% ---------------------------------------------------------------
% 1b. Primera iteración sobre la curva (n = 5, intervalo [0,100])
% ---------------------------------------------------------------
\begin{frame}{Fibonacci Search: Explicación Gráfica (2/3)}
    \begin{center}
    \begin{tikzpicture}[x=0.1cm, y=1cm, font=\footnotesize]
        \useasboundingbox (-4,-3.75) rectangle (104,2.4);
        % Función del ejemplo
        \node[anchor=north west] at (-2,2.35)
            {$f(x)=\sqrt{1+\left(\dfrac{x-45}{15}\right)^{2}}$,\quad $x\in[0,100]$};
        % f(x) = sqrt(1+((x-45)/15)^2), restada su mínimo para apoyarla en y = 0
        \draw[GrisTexto, thick, domain=0:100, samples=60, smooth, variable=\t]
            plot ({\t}, {(sqrt(1+((\t-45)/15)^2)-1)*0.6});
        \pgfmathsetmacro{\hc}{(sqrt(1+((37.5-45)/15)^2)-1)*0.6}
        \pgfmathsetmacro{\hd}{(sqrt(1+((62.5-45)/15)^2)-1)*0.6}

        % Paso 1: intervalo inicial [a,b]
        \draw[CelesteBrillante, very thick] (0,-0.45) -- (100,-0.45);
        \draw[CelesteBrillante, thick] (0,-0.55) -- (0,-0.35);
        \draw[CelesteBrillante, thick] (100,-0.55) -- (100,-0.35);
        \node[below] at (0,-0.55)   {$a=0$};
        \node[below] at (100,-0.55) {$b=100$};

        % Paso 2: puntos de consulta
        \foreach \px/\lb in {37.5/c, 62.5/d}{
            \draw[CelesteBrillante, thick] (\px,-0.55) -- (\px,-0.35);
            \node[below] at (\px,-0.55) {$\lb=\px$};
        }
        % Paso 3: se evalúa f en c y d
        \foreach \px/\py in {37.5/\hc, 62.5/\hd}{
            \draw[NaranjaBase, dashed] (\px,-0.35) -- (\px,\py);
            \node[circle, fill=NaranjaBase, inner sep=1.6pt] at (\px,\py) {};
        }
        \node[anchor=north east, text=NaranjaBase] at (36,\hc-0.03) {$f(c)$};
        \node[anchor=south east, text=NaranjaBase] at (61.5,\hd+0.03) {$f(d)$};
        % Paso 4: se descarta [d,b]; c queda como nuevo punto d
        \fill[RojoAlerta, opacity=0.25] (62.5,-0.35) rectangle (100,1.75);
        \draw[RojoAlerta, thick, dashed] (62.5,-1.45) -- (100,-1.45);
        \draw[CelesteBrillante, very thick] (0,-1.45) -- (62.5,-1.45);
        \draw[CelesteBrillante, thick] (0,-1.55) -- (0,-1.35);
        \draw[CelesteBrillante, thick] (62.5,-1.55) -- (62.5,-1.35);
        \node[circle, fill=VerdeNeon, inner sep=1.6pt] at (37.5,-1.45) {};
        \node[below, text=VerdeNeon] at (37.5,-1.55) {$c$ se reutiliza};

        % Panel de cálculos (dos columnas)
        \node[anchor=north west, text width=4.8cm, align=left] at (-2,-2.2)
            {$n=5,\quad \rho=F_5/F_6=5/8$};
        \node[anchor=north west, text width=4.8cm, align=left] at (-2,-2.7)
            {$d=(1-\rho)\,a+\rho\, b=62.5$};
        \node[anchor=north west, text width=4.8cm, align=left] at (-2,-3.2)
            {$c=\rho\, a+(1-\rho)\, b=37.5$};
        \node[anchor=north west, text width=4.8cm, align=left] at (52,-2.2)
            {$f(c)=1.12 < f(d)=1.54$};
        \node[anchor=north west, text width=4.8cm, align=left, text=RojoAlerta] at (52,-2.7)
            {Se descarta $[d,b]$:\\[2pt] $I_2=62.5=\tfrac{5}{8}\,I_1$};
    \end{tikzpicture}
    \end{center}
\end{frame}

% ---------------------------------------------------------------
% 1c. Evolución del intervalo (n = 5, intervalo [0,100])
% ---------------------------------------------------------------
% Una fila del diagrama: #1 altura | #2,#3 intervalo anterior | #4,#5 intervalo actual
%                        #6,#7 puntos izquierdo y derecho | #8,#9 sus colores
\newcommand{\fibrow}[9]{%
    \draw[gray!45, thin] (0,#1) -- (100,#1);
    \draw[RojoAlerta, thick, dashed] (#2,#1) -- (#4,#1);
    \draw[RojoAlerta, thick, dashed] (#5,#1) -- (#3,#1);
    \draw[CelesteBrillante, very thick] (#4,#1) -- (#5,#1);
    \foreach \e in {#4,#5}{\draw[CelesteBrillante, thick] (\e,{#1-0.09}) -- (\e,{#1+0.09});}
    \node[circle, fill=#8, inner sep=1.4pt] at (#6,#1) {};
    \node[circle, fill=#9, inner sep=1.4pt] at (#7,#1) {};
    \node[anchor=south east, inner sep=1pt] at (#6,{#1+0.06}) {$\pgfmathprintnumber[fixed,precision=1]{#6}$};
    \node[anchor=south west, inner sep=1pt] at (#7,{#1+0.06}) {$\pgfmathprintnumber[fixed,precision=1]{#7}$};
}

\begin{frame}{Fibonacci Search: Explicación Gráfica (3/3)}
    \begin{center}
    \begin{tikzpicture}[x=0.072cm, y=1cm, font=\scriptsize]
        \useasboundingbox (-3,-5.55) rectangle (146,0.5);
        \fibrow{0}{0}{100}{0}{100}{37.5}{62.5}{NaranjaBase}{NaranjaBase}
        \node[right, align=left] at (104,0) {$I_1=100=8u$\\ $\rho=5/8$};
        \fibrow{-1}{0}{100}{0}{62.5}{25}{37.5}{NaranjaBase}{VerdeNeon}
        \node[right, align=left] at (104,-1) {$I_2=62.5=5u$\\ $\rho=3/5$};
        \fibrow{-2}{0}{62.5}{25}{62.5}{37.5}{50}{VerdeNeon}{NaranjaBase}
        \node[right, align=left] at (104,-2) {$I_3=37.5=3u$\\ $\rho=2/3$};
        \fibrow{-3}{25}{62.5}{37.5}{62.5}{50}{50.125}{VerdeNeon}{NaranjaBase}
        \node[right, align=left] at (104,-3) {$I_4=25=2u$\\ $\rho=1/2$};
        \draw[gray!45, thin] (0,-4) -- (100,-4);
        \draw[RojoAlerta, thick, dashed] (50.125,-4) -- (62.5,-4);
        \draw[CelesteBrillante, very thick] (37.5,-4) -- (50.125,-4);
        \draw[CelesteBrillante, thick] (37.5,-4.09) -- (37.5,-3.91);
        \draw[CelesteBrillante, thick] (50.125,-4.09) -- (50.125,-3.91);
        \node[anchor=south east, inner sep=1pt] at (37.5,-3.94) {$37.5$};
        \node[anchor=south west, inner sep=1pt] at (50.125,-3.94) {$50.1$};
        \node[right] at (104,-4) {$I_5\approx 12.6\approx u$};
        \node[below] at (75,-4.15) {$u=100/F_6=12.5$};
        % Leyenda
        \draw[CelesteBrillante, very thick] (0,-5.2) -- (8,-5.2);
        \node[right] at (10,-5.2) {intervalo activo};
        \draw[RojoAlerta, thick, dashed] (44,-5.2) -- (52,-5.2);
        \node[right] at (54,-5.2) {descartado};
        \node[circle, fill=NaranjaBase, inner sep=1.4pt] at (82,-5.2) {};
        \node[right] at (85,-5.2) {punto nuevo};
        \node[circle, fill=VerdeNeon, inner sep=1.4pt] at (114,-5.2) {};
        \node[right] at (117,-5.2) {reutilizado};
    \end{tikzpicture}
    \end{center}
\end{frame}

% ---------------------------------------------------------------
% 2. Pseudocódigo
% ---------------------------------------------------------------
\begin{frame}{Fibonacci Search: Pseudocódigo}
    \scriptsize
    \begin{algorithmic}[1]
        \Require $f$ unimodal, intervalo $[a,b]$, presupuesto $n>1$, tolerancia $\epsilon$
        \State $\varphi \gets (1+\sqrt{5})/2$, \quad $s \gets (1-\sqrt{5})/(1+\sqrt{5})$
        \State $\rho \gets \left[\varphi\,(1-s^{n+1})/(1-s^{n})\right]^{-1}$
              \Comment{equivale a $F_n/F_{n+1}$}
        \State $d \gets (1-\rho)\,a + \rho\, b$; \quad $y_d \gets f(d)$; \quad $n \gets n-1$
        \While{$n > 0$}
            \If{$n > 1$}
                \State $c \gets \rho\, a + (1-\rho)\, b$ \Comment{punto simétrico de $d$}
            \Else
                \State $c \gets \epsilon\, a + (1-\epsilon)\, d$ \Comment{última consulta, pegada a $d$}
            \EndIf
            \State $\rho \gets \left[\varphi\,(1-s^{n+1})/(1-s^{n})\right]^{-1}$
            \State $y_c \gets f(c)$; \quad $n \gets n-1$
            \If{$y_c < y_d$}
                \State $(d,\, b,\, y_d) \gets (c,\, d,\, y_c)$ \Comment{el mínimo está del lado de $c$}
            \Else
                \State $(a,\, b) \gets (b,\, c)$ \Comment{se invierte la orientación}
            \EndIf
        \EndWhile
        \State \Return $[\min(a,b),\ \max(a,b)]$
    \end{algorithmic}
    \vspace{0.3em}
    {\tiny \textit{Nota:} el intercambio $(a,b)\gets(b,c)$ voltea el intervalo; por eso se ordena al retornar.}
\end{frame}

% ---------------------------------------------------------------
% 3. Ejemplo paso a paso
% ---------------------------------------------------------------
\begin{frame}{Fibonacci Search: Ejemplo Paso a Paso (1/3)}
    \scriptsize
    \begin{exampleblock}{Planteamiento}
        Minimizar $f(x)=\sqrt{1+\left(\frac{x-45}{15}\right)^{2}}$ en $[a,b]=[0,100]$
        con $n=5$ evaluaciones y $\epsilon=0.01$. El mínimo exacto es $x^{*}=45$.
    \end{exampleblock}

    \begin{block}{Inicialización (líneas 1--3)}
        $\varphi=1.618,\quad s=-0.382,\quad \rho=F_5/F_6=5/8=0.625$\\[2pt]
        $d=(1-\rho)\,a+\rho\, b=0.375\cdot 0+0.625\cdot 100=62.5$\\[2pt]
        $y_d=f(62.5)=1.5366,\quad n\gets 4$
    \end{block}

    \begin{block}{Iteración 1: $n=4>1$ (líneas 5--15)}
        $c=\rho\, a+(1-\rho)\, b=0.625\cdot 0+0.375\cdot 100=37.5,\quad \rho\gets F_4/F_5=3/5$\\[2pt]
        $y_c=f(37.5)=1.1180,\quad n\gets 3$\\[2pt]
        $y_c<y_d \;\Rightarrow\; (d,b,y_d)\gets(37.5,\;62.5,\;1.1180)$\\[2pt]
        Estado: $a=0,\; b=62.5,\; d=37.5 \;\Rightarrow\; I=[0,\,62.5]$
    \end{block}
\end{frame}

\begin{frame}{Fibonacci Search: Ejemplo Paso a Paso (2/3)}
    \scriptsize
    \begin{block}{Iteración 2: $n=3>1$}
        $c=\rho\, a+(1-\rho)\, b=0.6\cdot 0+0.4\cdot 62.5=25,\quad \rho\gets F_3/F_4=2/3$\\[2pt]
        $y_c=f(25)=1.6667,\quad n\gets 2$\\[2pt]
        $y_c\ge y_d\;(1.6667\ge 1.1180) \;\Rightarrow\; (a,b)\gets(b,c)=(62.5,\;25)$\\[2pt]
        Estado: $a=62.5,\; b=25,\; d=37.5 \;\Rightarrow\; I=[25,\,62.5]$\\[2pt]
        {\tiny Ahora $a>b$: el intervalo quedó volteado y se ordena recién al retornar (línea 18).}
    \end{block}

    \begin{block}{Iteración 3: $n=2>1$}
        $c=\tfrac{2}{3}\cdot 62.5+\tfrac{1}{3}\cdot 25=50,\quad \rho\gets F_2/F_3=1/2$\\[2pt]
        $y_c=f(50)=1.0541,\quad n\gets 1$\\[2pt]
        $y_c<y_d\;(1.0541< 1.1180) \;\Rightarrow\; (d,b,y_d)\gets(50,\;37.5,\;1.0541)$\\[2pt]
        Estado: $a=62.5,\; b=37.5,\; d=50 \;\Rightarrow\; I=[37.5,\,62.5]$
    \end{block}
\end{frame}

\begin{frame}{Fibonacci Search: Ejemplo Paso a Paso (3/3)}
    \scriptsize
    \begin{block}{Iteración 4: $n=1$, rama \textit{else} (línea 8)}
        $c=\epsilon\, a+(1-\epsilon)\, d=0.01\cdot 62.5+0.99\cdot 50=50.125$\\[2pt]
        $y_c=f(50.125)=1.0568,\quad n\gets 0$\\[2pt]
        $y_c\ge y_d\;(1.0568\ge 1.0541) \;\Rightarrow\; (a,b)\gets(b,c)=(37.5,\;50.125)$
    \end{block}

    \begin{exampleblock}{Resultado (línea 18)}
        $[\min(a,b),\,\max(a,b)]=[37.5,\;50.125]$, de longitud $12.625\approx 100/F_6=12.5$,
        que contiene a $x^{*}=45$.
    \end{exampleblock}

    \begin{center}
    \begin{tabular}{l c c c c c}
        \toprule
        & inicio & tras it.\,1 & tras it.\,2 & tras it.\,3 & tras it.\,4 \\
        \midrule
        $[a,b]$  & $[0,100]$ & $[0,62.5]$ & $[25,62.5]$ & $[37.5,62.5]$ & $[37.5,50.125]$ \\
        longitud & $100$ & $62.5$ & $37.5$ & $25$ & $12.625$ \\
        múltiplo de $u=12.5$ & $8u$ & $5u$ & $3u$ & $2u$ & $\approx u$ \\
        \bottomrule
    \end{tabular}
    \end{center}
\end{frame}

% ---------------------------------------------------------------
% 4. Código
% ---------------------------------------------------------------
\begin{frame}[fragile]{Fibonacci Search: Código}
\begin{lstlisting}[language=Python]
from math import sqrt
PHI = (1 + sqrt(5)) / 2
def fibonacci_search(f, a, b, n, eps=0.01):
    s = (1 - sqrt(5)) / (1 + sqrt(5))
    rho = 1 / (PHI * (1 - s**(n + 1)) / (1 - s**n))
    d = (1 - rho) * a + rho * b       # punto derecho
    yd, n = f(d), n - 1
    while n > 0:
        if n > 1:
            c = rho * a + (1 - rho) * b   # punto izquierdo
        else:
            c = eps * a + (1 - eps) * d   # ultima consulta
        rho = 1 / (PHI * (1 - s**(n + 1)) / (1 - s**n))
        yc, n = f(c), n - 1
        if yc < yd:
            d, b, yd = c, d, yc   # el minimo esta junto a c
        else:
            a, b = b, c           # se invierte el intervalo
    return (a, b) if a < b else (b, a)
# f(x) = sqrt(1 + ((x-45)/15)**2)
# fibonacci_search(f, 0, 100, 5) -> (37.5, 50.125)
\end{lstlisting}
\end{frame}

% ---------------------------------------------------------------
% 5. Análisis de convergencia
% ---------------------------------------------------------------
\begin{frame}{Fibonacci Search: Análisis de Convergencia}
    \begin{alertblock}{Tasa de convergencia}
        \small
        Tras $n$ evaluaciones la incertidumbre es
        $\;|b_n-a_n| = \dfrac{b-a}{F_{n+1}}$.
        Como $F_{n+1}\approx \varphi^{\,n+1}/\sqrt{5}$, el decaimiento es exponencial
        y la convergencia es \textbf{lineal} con razón asintótica
        $1/\varphi \approx 0.618$.
    \end{alertblock}

    \footnotesize
    \begin{itemize}
        \item \textbf{Orden $q=1$:} $\varepsilon_{k+1}/\varepsilon_k = F_{k}/F_{k+1} \to 1/\varphi$;
              no hay aceleración superlineal porque solo se usa información de orden cero.
        \item \textbf{Optimalidad:} ningún método basado en comparaciones logra, con $n$
              consultas, un intervalo garantizado más pequeño que $(b-a)/F_{n+1}$.
        \item \textbf{Costo:} $1$ evaluación por iteración y sin derivadas; para una
              tolerancia $\epsilon$ basta el menor $n$ con $F_{n+1}\ge (b-a)/\epsilon$.
        \item \textbf{Limitación:} $n$ debe fijarse de antemano, pues $\rho$ depende de
              él; si $f$ no es unimodal, solo se garantiza un mínimo local.
    \end{itemize}
\end{frame}
